\documentclass[12pt]{article}
\usepackage[utf8]{inputenc}
\usepackage[T1]{fontenc}
\usepackage[margin=1in]{geometry}
\usepackage{setspace}
\usepackage{xcolor}
\usepackage[colorlinks=true, urlcolor=blue!60!black, linkcolor=blue!60!black]{hyperref}
\setlength{\parindent}{2em}
\setlength{\parskip}{1.25em}
\renewcommand{\baselinestretch}{1.0}
\begin{document}
\begin{center}
{\LARGE \textit{Learning Resources, Inc. v. Trump}, 607 U.S. \_\_\_ (2026)\par}
\vspace{0.8em}
{\large Case Brief\par}
\end{center}
\vspace{1em}

\section{Memory Jogger}
In \href{https://www.courtlistener.com/opinion/10799354/learning-resources-inc-v-trump-revisions-22326/}{\textit{Learning Resources, Inc. v. Trump}}, 607 U.S. \_\_\_ (2026) (No. 24-1287, consolidated with \textit{Trump v. V.O.S. Selections, Inc.}, No. 25-250), the Supreme Court held 6--3 that the International Emergency Economic Powers Act's grant of authority to ``regulate \ldots{} importation'' does not authorize the President to impose tariffs, while fracturing sharply over the foundation and scope of the Major Questions Doctrine and the Nondelegation Doctrine.

\section{Detailed Case Facts}
Shortly after taking office in early 2025, President Donald J. Trump declared national emergencies under the National Emergencies Act (NEA), 50 U.S.C. \S~1601 \textit{et seq.}, and invoked the International Emergency Economic Powers Act (IEEPA), 50 U.S.C. \S\S~1701--1702, to address two perceived foreign threats (\href{https://www.law.cornell.edu/supremecourt/text/24-1287}{607 U.S. \_\_\_, slip op. at 1--2}). Enacted in 1977, IEEPA authorizes the President, upon declaring a national emergency to deal with an ``unusual and extraordinary threat'' originating primarily outside the United States, to ``investigate, block during the pendency of an investigation, regulate, direct and compel, nullify, void, prevent or prohibit \ldots{} importation or exportation of \ldots{} any property in which any foreign country or a national thereof has any interest'' (50 U.S.C. \S\S~1701(a), 1702(a)(1)(B); \href{https://www.law.cornell.edu/supremecourt/text/24-1287}{slip op. at 2}).

First, targeting the influx of illegal drugs (particularly fentanyl and opioids) from Canada, Mexico, and China, President Trump issued Presidential Proclamation No. 10886 and Executive Orders Nos. 14193, 14194, and 14195 (90 Fed. Reg. 8327, 9113, 9117, 9121 (2025)), imposing ``drug trafficking'' tariffs of 25\% on most Canadian and Mexican imports and 10\% on most Chinese imports (\href{https://www.law.cornell.edu/supremecourt/text/24-1287}{slip op. at 1--2}). He subsequently doubled the Chinese drug-trafficking tariff to 20\% (Exec. Order No. 14228) and eliminated the statutory \$800 \textit{de minimis} exemption for Chinese goods (Exec. Order No. 14256) (\href{https://www.law.cornell.edu/supremecourt/text/24-1287}{slip op. at 3}).

Second, citing ``large and persistent'' trade deficits that he determined had hollowed out American manufacturing and undermined critical supply chains, President Trump issued Executive Order No. 14257 (90 Fed. Reg. 15041 (2025)), imposing worldwide ``reciprocal'' tariffs of at least 10\% on all imports from all trading partners, with higher country-specific rates for dozens of nations, notwithstanding existing trade agreements (\href{https://www.law.cornell.edu/supremecourt/text/24-1287}{slip op. at 1--3}). Through rapid follow-on orders (Exec. Orders Nos. 14259 and 14266), he escalated reciprocal tariffs on Chinese goods from 34\% to 84\% and then 125\%, bringing the total effective tariff rate on most Chinese goods to 145\%, while repeatedly modifying rates and product exemptions (e.g., Exec. Orders Nos. 14346, 14358, and 14360) (\href{https://www.law.cornell.edu/supremecourt/text/24-1287}{slip op. at 3}). Over the months the tariffs remained in effect, the federal government collected more than \$200 billion from importers (\href{https://www.theregreview.org/2026/07/24/levin-the-tariff-case-and-the-major-questions-doctrine/}{The Regulatory Review}).

Petitioners Learning Resources, Inc. and hand2mind, Inc.---two family-owned businesses designing and distributing children's educational products largely manufactured abroad, including in China---faced an existential spike in import duties from \$2.3 million in 2024 to over \$100 million projected for 2025 (\href{https://www.oyez.org/cases/2025/24-1287}{Oyez}). Meanwhile, in a parallel suit (\textit{V.O.S. Selections, Inc. v. Trump}), five small businesses and twelve States likewise challenged both sets of IEEPA tariffs (\href{https://www.law.cornell.edu/supremecourt/text/24-1287}{slip op. at 3--4}).

\section{Procedural History}
On April 22, 2025, Learning Resources, Inc. and hand2mind, Inc. filed suit in the U.S. District Court for the District of Columbia challenging the IEEPA tariffs (\href{https://www.oyez.org/cases/2025/24-1287}{Oyez}). The Government moved to transfer the case to the U.S. Court of International Trade (CIT) under 28 U.S.C. \S~1581(i)(1), which grants the CIT exclusive jurisdiction over civil actions against the United States arising out of laws providing for tariffs or their administration (\href{https://www.law.cornell.edu/supremecourt/text/24-1287}{slip op. at 4}). On May 29, 2025, District Judge Rudolph Contreras denied the motion to transfer and granted a preliminary injunction, holding that IEEPA does not authorize the President to impose tariffs (\textit{Learning Resources, Inc. v. Trump}, 784 F. Supp. 3d 209 (D.D.C. 2025); \href{https://www.law.cornell.edu/supremecourt/text/24-1287}{slip op. at 4}). The district court stayed its injunction pending appeal to the D.C. Circuit, and petitioners sought certiorari before judgment (\href{https://www.oyez.org/cases/2025/24-1287}{Oyez}).

Separately, on May 28, 2025, a three-judge panel of the CIT granted summary judgment to the \textit{V.O.S. Selections} plaintiffs (five small businesses and twelve States) (\textit{V.O.S. Selections, Inc. v. United States}, 772 F. Supp. 3d 1350 (Ct. Int'l Trade 2025); \href{https://www.law.cornell.edu/supremecourt/text/24-1287}{slip op. at 4}). After staying the CIT's judgment, the U.S. Court of Appeals for the Federal Circuit, sitting \textit{en banc}, affirmed in relevant part on August 29, 2025 (149 F. 4th 1312 (Fed. Cir. 2025)) (\href{https://www.law.cornell.edu/supremecourt/text/24-1287}{slip op. at 4}). The Federal Circuit held that the CIT possessed exclusive jurisdiction because the claims arose out of modifications to the Harmonized Tariff Schedule of the United States (HTSUS), and on the merits agreed that IEEPA's authorization to ``regulate \ldots{} importation'' did not authorize tariffs ``unbounded in scope, amount, and duration'' (Judge Cunningham concurred for four judges that IEEPA authorizes no tariffs at all; Judge Taranto dissented for four judges) (\href{https://www.law.cornell.edu/supremecourt/text/24-1287}{slip op. at 4}).

On September 9, 2025, the Supreme Court granted the Government's expedited petition for certiorari in \textit{V.O.S. Selections} (No. 25-250) and the \textit{Learning Resources} petition for certiorari before judgment (No. 24-1287), consolidating both cases (606 U.S. 1050 (2025); \href{https://www.scotusblog.com/cases/learning-resources-inc-v-trump/}{SCOTUSblog}). Oral argument took place on November 5, 2025, and the Court issued its decision on February 20, 2026 (\href{https://www.supremecourt.gov/opinions/25pdf/24-1287_4gcj.pdf}{Official Slip Opinion}).

\section{Judicial Votes}
The Supreme Court ruled 6--3 against the President's IEEPA tariffs, producing seven separate opinions (\href{https://www.scotusblog.com/cases/learning-resources-inc-v-trump/}{SCOTUSblog}; \href{https://www.law.cornell.edu/supremecourt/text/24-1287}{Cornell LII}):
\begin{itemize}
\item \textbf{Opinion of the Court (Parts I, II--A--1, and II--B) (6 Justices):} Chief Justice John G. Roberts, Jr. delivered the opinion of the Court as to Parts I, II--A--1, and II--B, joined by Justices Sonia Sotomayor, Elena Kagan, Neil M. Gorsuch, Amy Coney Barrett, and Ketanji Brown Jackson.
\item \textbf{Plurality Opinion (Parts II--A--2 and III) (3 Justices):} Chief Justice Roberts authored Parts II--A--2 (invoking the Major Questions Doctrine) and III, joined only by Justices Gorsuch and Barrett.
\item \textbf{Concurrences Joining the Principal Opinion in Full (2 Justices):} Justice Gorsuch and Justice Barrett each filed separate concurring opinions debating the doctrinal nature of the Major Questions Doctrine.
\item \textbf{Concurrences in Part and in the Judgment (3 Justices):} Justice Kagan filed an opinion concurring in part and concurring in the judgment, joined by Justices Sotomayor and Jackson (rejecting reliance on the Major Questions Doctrine). Justice Jackson also filed a separate opinion concurring in part and concurring in the judgment (relying on legislative history).
\item \textbf{Dissents (3 Justices):} Justice Brett M. Kavanaugh filed the principal dissenting opinion, joined by Justices Clarence Thomas and Samuel A. Alito. Justice Thomas also filed a separate dissenting opinion on originalist nondelegation grounds.
\end{itemize}

\section{Holding}
\textbf{Judgment:} In No. 25-250 (\textit{Trump v. V.O.S. Selections, Inc.}), the judgment of the U.S. Court of Appeals for the Federal Circuit was \textbf{affirmed}. In No. 24-1287 (\textit{Learning Resources, Inc. v. Trump}), the judgment of the U.S. District Court for the District of Columbia was \textbf{vacated, and the case was remanded with instructions to dismiss for lack of jurisdiction} because the CIT has exclusive jurisdiction under 28 U.S.C. \S~1581(i)(1) (\href{https://www.law.cornell.edu/supremecourt/text/24-1287}{slip op. at 4 n.2, 21}).

\textbf{Rule Announced:} The International Emergency Economic Powers Act, 50 U.S.C. \S~1702(a)(1)(B), does not authorize the President to impose tariffs. Because Article I, Section 8 vests the taxing and tariff power exclusively in Congress and the President possesses no inherent peacetime tariff authority, IEEPA's grant of authority to ``regulate \ldots{} importation''---read in textual, structural, and historical context---does not encompass the distinct power to levy revenue-raising customs duties (\href{https://www.supremecourt.gov/opinions/25pdf/24-1287_4gcj.pdf}{slip op. at 5--7, 14--20}).

\textbf{Holding vs. Dicta:} The binding six-Justice holding rests strictly on ordinary statutory interpretation in Parts I, II--A--1, and II--B: IEEPA does not authorize tariffs, and challenges to HTSUS-modifying tariff orders fall within the CIT's exclusive jurisdiction. Crucially, Chief Justice Roberts's Major Questions Doctrine analysis in Part II--A--2 (including his statements that there is no emergency or foreign-affairs exception to the Major Questions Doctrine) commanded only a three-Justice plurality (Roberts, Gorsuch, and Barrett) and is therefore \textbf{not} a binding majority holding. Moreover, the Court expressly disclaimed setting ``the metes and bounds'' of non-tariff authority to ``regulate \ldots{} importation'' under IEEPA, declined to speculate on whether the President could reimpose similar tariffs under Title 19 statutes (such as Section 232 or Section 301), did not decide whether presidential emergency declarations under IEEPA are judicially reviewable, and did not address how refunds of collected IEEPA tariffs must be administered (\href{https://www.law.cornell.edu/supremecourt/text/24-1287}{slip op. at 16 \& n.5}; \href{https://www.lawfaremedia.org/article/constitutional-duels-in-the-court-s-rejection-of-trump-s-tariffs}{Lawfare}).

\section{Analysis of Opinions}

\subsection{Opinion of the Court (Parts I, II--A--1, and II--B) and Plurality Opinion (Parts II--A--2 and III) (Chief Justice Roberts)}
Writing for a six-Justice majority in Parts I and II--A--1, Chief Justice Roberts grounded the analysis in Article I, Section 8, Clause 1 (granting Congress the power to ``lay and collect Taxes, Duties, Imposts and Excises'') and the Origination Clause (Art. I, \S~7, cl. 1), emphasizing that tariffs are ``very clear[ly] \ldots{} a branch of the taxing power'' (\textit{Gibbons v. Ogden}) and that the Framers gave ``Congress alone \ldots{} access to the pockets of the people'' (\href{https://www.law.cornell.edu/supremecourt/text/24-1287}{slip op. at 5--6}). Because the Government conceded the President has no inherent peacetime authority to impose tariffs and did not invoke war powers, its defense stood or fell solely on IEEPA (\href{https://www.law.cornell.edu/supremecourt/text/24-1287}{slip op. at 6--7}).

In Part II--A--2 (joined only by Justices Gorsuch and Barrett), Roberts invoked the Major Questions Doctrine (\textit{West Virginia v. EPA}; \textit{Biden v. Nebraska}; \textit{NFIB v. OSHA}). He stressed that Congress consistently delegates tariff authority expressly and with strict rate caps, time limits, and procedural prerequisites in Title 19; that under the NEA (post-\textit{INS v. Chadha}), Congress can terminate a presidential emergency only via a joint resolution overcoming a presidential veto (50 U.S.C. \S~1622(a)(1)); that no President in IEEPA's nearly 50-year history had ever used it to impose tariffs; and that the economic stakes (\$4 trillion in projected deficit reduction and \$15 trillion in trade deals) dwarfed prior major-questions cases (\href{https://www.law.cornell.edu/supremecourt/text/24-1287}{slip op. at 7--11}). Rejecting proposed emergency and foreign-affairs exceptions, Roberts declared that ``[t]here is no major questions exception to the major questions doctrine'' (\href{https://www.law.cornell.edu/supremecourt/text/24-1287}{slip op. at 11--13}).

Returning to a six-Justice majority in Part II--B, Roberts held that IEEPA's text and structure do not confer tariff power: (1) the statute omits any mention of ``tariffs'' or ``duties''; (2) the ordinary meaning of ``regulate'' does not include taxation, and the Government could not identify a single statute in the U.S. Code where ``regulate'' authorizes a tax; (3) reading ``regulate'' to include taxation would render IEEPA's parallel authorization to ``regulate \ldots{} exportation'' unconstitutional under the Export Clause (Art. I, \S~9, cl. 5); (4) under \textit{noscitur a sociis}, the nine verbs in \S~1702(a)(1)(B) authorize sanctions and transaction controls, not revenue-raising taxes; and (5) neither the greater-includes-the-lesser argument, nor \textit{United States v. Yoshida International, Inc.}, 526 F.2d 560 (CCPA 1975) (a single intermediate court decision under IEEPA's predecessor, the Trading with the Enemy Act (TWEA)), nor \textit{FEA v. Algonquin SNG, Inc.}, 426 U.S. 548 (1976), nor \textit{Dames \& Moore v. Regan}, 453 U.S. 654 (1981), supported reading tariffs into IEEPA (\href{https://www.law.cornell.edu/supremecourt/text/24-1287}{slip op. at 14--20}).

\subsection{Concurring Opinion (Justice Gorsuch)}
Justice Gorsuch joined Roberts's opinion in full and wrote an expansive solo concurrence defending the Major Questions Doctrine as a Constitution-enforcing substantive clear-statement rule rooted in Article I's Legislative Vesting Clause and centuries of common-law corporate, agency, and executive-power clear-statement rules (\href{https://www.law.cornell.edu/supremecourt/text/24-1287}{Gorsuch, J., concurring, slip op. at 1--17}). Highlighting the post-\textit{Chadha} retrieval problem---once a court reads a vague statute to hand power to the President, Congress can retrieve it only by overcoming a presidential veto---Gorsuch argued the doctrine is ``pro-Congress,'' not ``anti-administrative state'' (\href{https://www.law.cornell.edu/supremecourt/text/24-1287}{slip op. at 15--17}). He critiqued all three other methodological camps: (1) chiding Justice Kagan for ``all but apply[ing]'' the Major Questions Doctrine here while reading broad statutes expansively in \textit{West Virginia} and \textit{Nebraska}; (2) contending against Justice Barrett that ``commonsense principles of communication'' cannot explain major-questions cases without relying on external Article I norms; (3) rejecting Justice Kavanaugh's foreign-affairs exception except where independent Article II powers overlap with Congress's; and (4) refuting Justice Thomas's narrow reading of the nondelegation doctrine as contrary to Article I's text, Founding debates, and the Revolutionary rejection of executive duties (\href{https://www.law.cornell.edu/supremecourt/text/24-1287}{slip op. at 2--46}).

\subsection{Concurring Opinion (Justice Barrett)}
Justice Barrett joined the Chief Justice's opinion in full and wrote briefly to respond to Justice Gorsuch (\href{https://www.law.cornell.edu/supremecourt/text/24-1287}{Barrett, J., concurring, slip op. at 1}). Reaffirming her concurrence in \textit{Biden v. Nebraska}, she insisted that the Major Questions Doctrine is not a ``strong-form substantive canon'' imposing a ``clarity tax'' or ``judicial flex'' that forces judges to abandon the most natural reading of a statute (\href{https://www.law.cornell.edu/supremecourt/text/24-1287}{slip op. at 1--4}). Rather, it is an ordinary application of contextual textualism that reads statutory text alongside background legal conventions, common sense, and constitutional structure (including Article I) to determine a statute's most natural meaning (\href{https://www.law.cornell.edu/supremecourt/text/24-1287}{slip op. at 1--2}). Because strong-form dice-loading canons veer from interpretation into judicial policymaking, Barrett rejected Gorsuch's substantive-canon project as in tension with textualism (\href{https://www.law.cornell.edu/supremecourt/text/24-1287}{slip op. at 3--4}).

\subsection{Opinion Concurring in Part and Concurring in the Judgment (Justice Kagan)}
Justice Kagan, joined by Justices Sotomayor and Jackson, joined Parts I, II--A--1, and II--B of Roberts's opinion but declined to join Parts II--A--2 and III because ``straight-up statutory construction'' resolved the case without any ``major-questions thumb on the interpretive scales'' (\href{https://www.law.cornell.edu/supremecourt/text/24-1287}{Kagan, J., concurring in part, slip op. at 1--3, 7}). Combining IEEPA's 9 verbs and 11 objects yields 99 authorized actions, none of the other 98 of which involves raising revenue; reading ``regulate \ldots{} importation'' to permit unlimited tariffs would make it the ``odd man out'' and gut Congress's carefully constrained tariff statutes in Title 19 (\href{https://www.law.cornell.edu/supremecourt/text/24-1287}{slip op. at 4--6}). Answering Justice Gorsuch in a pointed footnote, Kagan wrote that ``I almost regret to inform him that I am not'' a convert to the Major Questions Doctrine, explaining that ordinary text-in-context reading naturally yields different results across differently structured delegation statutes (\href{https://www.law.cornell.edu/supremecourt/text/24-1287}{slip op. at 3 n.1}).

\subsection{Opinion Concurring in Part and Concurring in the Judgment (Justice Jackson)}
Justice Jackson joined Parts I, II--A--1, and II--B of Roberts's opinion and Justice Kagan's concurrence, writing separately to argue that the Court should consult legislative history rather than speculate about congressional intent through the Major Questions Doctrine (\href{https://www.law.cornell.edu/supremecourt/text/24-1287}{Jackson, J., concurring in part, slip op. at 1--2}). Examining the 1941 House and Senate Reports on the TWEA amendment and the 1977 House and Senate Reports on IEEPA, Jackson demonstrated that Congress intended the ``regulate \ldots{} importation'' provision solely to empower the President ``to control or freeze property transactions where a foreign interest is involved''---not to levy revenue-raising tariffs (\href{https://www.law.cornell.edu/supremecourt/text/24-1287}{slip op. at 2--5}).

\subsection{Principal Dissent (Justice Kavanaugh)}
Justice Kavanaugh, joined by Justices Thomas and Alito, argued that IEEPA's text, statutory history, and Supreme Court precedent authorize presidential tariffs (\href{https://www.naag.org/attorney-general-journal/supreme-court-report-volume-33-issue-7/}{NAAG Supreme Court Report}; \href{https://www.hunton.com/insights/legal/major-questions-remain-about-the-major-questions-doctrine}{Hunton Analysis}). Textually, Kavanaugh contended that tariffs are a traditional and common means to ``regulate importation''---posing the question that if quotas and embargoes are means to regulate importation, ``how are tariffs not a means to regulate importation?'' (\href{https://supreme.justia.com/cases/federal/us/607/24-1287/}{Justia}). Relying on President Nixon's 1971 surcharge upheld in \textit{Yoshida}, \textit{Algonquin}, and \textit{Dames \& Moore}, he maintained that Congress incorporated tariff authority when it carried TWEA's wording into IEEPA (\href{https://www.naag.org/attorney-general-journal/supreme-court-report-volume-33-issue-7/}{NAAG}).

On the Major Questions Doctrine, Kavanaugh argued both that IEEPA provides clear authorization (criticizing the plurality for imposing a ``magic-words test'') and that the Major Questions Doctrine does not apply to foreign-affairs statutes, where Congress traditionally grants the President broad flexibility (\href{https://www.hunton.com/insights/legal/major-questions-remain-about-the-major-questions-doctrine}{Hunton Analysis}). Finally, Kavanaugh observed that the President could likely reimpose ``most if not all'' of the tariffs under alternative Title 19 statutes (such as the Trade Expansion Act of 1962, the Trade Act of 1974, and the Tariff Act of 1930) and warned that refunding billions of dollars already collected ``is likely to be a `mess'\,'' (\href{https://www.law.cornell.edu/supremecourt/text/24-1287}{slip op. at 16 n.5}; \href{https://www.naag.org/attorney-general-journal/supreme-court-report-volume-33-issue-7/}{NAAG}).

\subsection{Separate Dissent (Justice Thomas)}
Justice Thomas joined Justice Kavanaugh's dissent in full and wrote separately to articulate an originalist theory of the Nondelegation Doctrine (\href{https://www.law.cornell.edu/supremecourt/text/24-1287}{Thomas, J., dissenting, slip op. at 1--2}). Rooted in both Article I's Legislative Vesting Clause and the Fifth Amendment's Due Process Clause, Thomas argued that the nondelegation bar applies \textit{only} to ``core legislative power''---making generally applicable rules setting conditions for deprivations of core private rights to life, liberty, or property (\href{https://www.law.cornell.edu/supremecourt/text/24-1287}{slip op. at 2--5}). Other powers assigned to Congress in Article I, Section 8---including foreign commerce and import duties---were historically royal prerogative powers governing privileges rather than private rights (\href{https://www.law.cornell.edu/supremecourt/text/24-1287}{slip op. at 5--12}). Because ``[i]mporting is a matter of privilege,'' Congress may delegate tariff power broadly to the President without implicating constitutional separation-of-powers limits (\href{https://www.law.cornell.edu/supremecourt/text/24-1287}{slip op. at 10--12}).

\section{Examples: Future Applications}
\begin{enumerate}
\item \textbf{Same-Side (Invalidated): Presidential Carbon Border Tax Under IEEPA.} Suppose a future President declares a national emergency over global climate change under the NEA and invokes IEEPA's authority to ``regulate \ldots{} importation'' to levy a 30\% border carbon tax on all imported gas-powered automobiles and steel (the exact hypothetical Justice Gorsuch highlighted at oral argument and in his concurrence). This falls squarely on the same side as \textit{Learning Resources}: it uses IEEPA's non-taxation verb ``regulate'' to raise revenue from domestic importers without express congressional tariff authorization.
\item \textbf{Same-Side (Invalidated): Executive Agency Fee or Tax Under a General ``Regulate'' Statute.} Suppose the Securities and Exchange Commission invokes its statutory authority to ``regulate the trading of \ldots{} securities'' (15 U.S.C. \S~78i(h)(1)) to impose a 1\% financial-transactions tax on stock trades to curb market speculation and raise Treasury revenue. Under \textit{Learning Resources}'s six-Justice textual holding, the quotidian statutory power to ``regulate'' does not encompass the distinct constitutional power to tax; Congress must confer taxing authority expressly.
\item \textbf{Opposite-Side (Upheld): IEEPA Asset Freezes and Targeted Import Embargoes Against Foreign Adversaries.} Suppose the President declares a national emergency over a foreign adversary's cyberattacks and invokes IEEPA to freeze that country's sovereign assets in U.S. banks and prohibit all imports of its telecommunications equipment. This falls squarely on the opposite side: \S~1702(a)(1)(B) expressly authorizes the President to ``block,'' ``prevent,'' or ``prohibit'' importation and foreign-property transactions, matching decades of presidential practice and IEEPA's core sanctions purpose without exercising the taxing power.
\item \textbf{Opposite-Side (Upheld): Sector-Specific Tariffs Imposed Under Title 19 After Statutory Procedures.} Suppose the President imposes a 15\% temporary import surcharge for 120 days to address a serious balance-of-payments deficit under Section 122 of the Trade Act of 1974 (19 U.S.C. \S~2132(a)), or imposes steel duties after a formal Commerce Department national-security investigation and report under Section 232 of the Trade Expansion Act of 1962 (19 U.S.C. \S~1862). This falls on the opposite side because Congress expressly delegated ``duty'' or ``surcharge'' authority in Title 19 subject to explicit rate caps, duration limits, and procedural prerequisites.
\item \textbf{Fence-Sitter: Presidential Licensing Fees on Foreign Software or Data Imports Under IEEPA, or Wartime Tariffs.} Suppose the President invokes IEEPA to impose a ``regulatory licensing fee'' (calibrated only to cover administrative oversight costs rather than general revenue) on foreign AI models imported into the United States, or imposes tariffs during an active congressionally authorized armed conflict invoking concurrent Commander-in-Chief powers. This is a fence-sitter: challengers would argue any monetary exaction on imports is a disguised tariff barred by \textit{Learning Resources}, whereas the Executive would distinguish cost-recovery regulatory fees from revenue-raising taxes and invoke Justice Gorsuch's and Chief Justice Roberts's acknowledgment that overlapping wartime or inherent Article II powers present a different constitutional calculus.
\end{enumerate}

\section{Critique}
\textit{Learning Resources, Inc. v. Trump} immediately drew intense scholarly scrutiny---not only for checking unilateral executive tariff power, but for exposing deep fissures within the Court's textualist majority over the Major Questions Doctrine and the Nondelegation Doctrine.

From a \textbf{textualist and originalist perspective}, commentators observed that \textit{Learning Resources} revealed an open ``canon war'' among conservative Justices. Yale Law Professor Abbe R. Gluck noted that all seven opinions grappled with the Major Questions Doctrine and exposed a sharp methodological rift between Justice Gorsuch---who candidly defends the doctrine as an external, Constitution-enforcing clear-statement rule even when statutory text is unambiguous---and Justice Barrett, who fears that strong-form substantive canons impose an atextual ``clarity tax'' and ``judicial flex'' on Congress (\href{https://www.scotusblog.com/2026/02/the-major-debate-over-major-questions-in-the-tariffs-decision-is-only-the-beginning/}{Gluck, SCOTUSblog (Feb. 27, 2026)}). Meanwhile, Justice Thomas's separate dissent drew originalist attention for attempting to narrow the Nondelegation Doctrine strictly to deprivations of ``life, liberty, or property,'' a move Gorsuch criticized as erasing Article I's text and inviting permanent executive aggrandizement over spending, borrowing, and foreign trade.

From a \textbf{progressive and administrative-law perspective}, scholars welcomed the Court's willingness to apply separation-of-powers limits evenhandedly to a Republican President's signature initiative, while questioning the ongoing coherence and necessity of the Major Questions Doctrine. Ronald M. Levin argued that \textit{Learning Resources} left the Major Questions Doctrine ``more tenuous'' because only three Justices joined the MQD portion of Roberts's opinion, Barrett retreated to ordinary contextual textualism, and Kagan demonstrated that standard statutory interpretation sufficed without a judge-made clear-statement hammer (\href{https://www.theregreview.org/2026/07/24/levin-the-tariff-case-and-the-major-questions-doctrine/}{Levin, The Regulatory Review (July 24, 2026)}). Michael R. Dreeben similarly emphasized that while the Court erected an important guardrail against unilateral emergency taxation, the ruling left unresolved key questions about wartime tariff loopholes, judicial reviewability of emergency declarations, and Congress's institutional willingness to reclaim powers ``slipping through its fingers'' after \textit{INS v. Chadha} eliminated the legislative veto (\href{https://www.lawfaremedia.org/article/constitutional-duels-in-the-court-s-rejection-of-trump-s-tariffs}{Dreeben, Lawfare (Mar. 11, 2026)}). Legal scholars also noted that \textit{Learning Resources} preserves bounded emergency governance under IEEPA for non-tariff sanctions while reinforcing that emergency statutes do not suspend Article I's allocation of the taxing power (\href{https://illinoislawreview.org/online/emergency-economic-powers-and-the-non-tariff-limits-of-ieepa/}{Illinois Law Review Online (2026)}).

At its core, \textit{Learning Resources} illuminates enduring \textbf{competing constitutional values}: democratic accountability, bicameral deliberation, and legislative control over the power of the purse (especially where post-\textit{Chadha} veto rules make delegated emergency powers nearly impossible for Congress to retrieve) against executive speed, energy, and negotiating flexibility in foreign economic crises; and judicial fidelity to enacted statutory text against judge-made clear-statement canons designed to enforce structural separation of powers.

\section{Key Quotations}
\begin{quote}
``Based on two words separated by 16 others in Section 1702(a)(1)(B) of IEEPA---`regulate' and `importation'---the President asserts the independent power to impose tariffs on imports from any country, of any product, at any rate, for any amount of time. Those words cannot bear such weight.''
\end{quote}
\noindent --- Chief Justice John G. Roberts, Jr., Opinion of the Court, \href{https://www.law.cornell.edu/supremecourt/text/24-1287}{607 U.S. \_\_\_, slip op. at 5 (2026)}.

\begin{quote}
``There is no major questions exception to the major questions doctrine.''
\end{quote}
\noindent --- Chief Justice John G. Roberts, Jr., Plurality Opinion, \href{https://www.law.cornell.edu/supremecourt/text/24-1287}{607 U.S. \_\_\_, slip op. at 13 (2026)}.

\begin{quote}
``The major questions doctrine is not `anti-administrative state.' It is pro-Congress.''
\end{quote}
\noindent --- Justice Neil M. Gorsuch, Concurring, \href{https://www.law.cornell.edu/supremecourt/text/24-1287}{607 U.S. \_\_\_, slip op. at 15 (2026)}.

\begin{quote}
``For all those reasons, straight-up statutory construction resolves this case for me; I need no major-questions thumb on the interpretive scales.''
\end{quote}
\noindent --- Justice Elena Kagan, Concurring in Part and Concurring in the Judgment, \href{https://www.law.cornell.edu/supremecourt/text/24-1287}{607 U.S. \_\_\_, slip op. at 7 (2026)}.

\begin{quote}
``Both Clauses forbid Congress from delegating core legislative power, which is the power to make substantive rules setting the conditions for deprivations of life, liberty, or property. Neither Clause prohibits Congress from delegating other kinds of power.''
\end{quote}
\noindent --- Justice Clarence Thomas, Dissenting, \href{https://www.law.cornell.edu/supremecourt/text/24-1287}{607 U.S. \_\_\_, slip op. at 2 (2026)}.


\end{document}
